\section{从诛吕到七国：恢复如何留下新的难题}
\label{sec:western-han-recovery-law-princes}

前一八〇年吕后去世后，周勃、陈平等诛灭吕氏，代王刘恒入长安即位，是为文帝。政变参与者的分工、宫廷内外如何联络以及少帝的处境，在《史记》《汉书》之间并非没有疑点；吕氏政权瓦解、功臣选择一位刘氏成年宗室继位的结果，则可以确证。它显示汉初政治并不只由皇帝血统维系：外戚、功臣、宗室和宫廷军力都须在继承时重新找到可接受的位置。

\subsection{减轻负担不等于没有国家}

文帝前一七八年诏减田租为三十税一，后来常被概括为“轻徭薄赋”。诏令能够证明朝廷试图以较低的田租稳定战后生产，却不能给出各地实际税率、附加劳役或农户所得的全国平均数。把一份中央命令直接等同于百姓共同经验，会抹去郡县执行、土地占有和家庭处境的差异。它仍是一项重要选择：中央需要重新获得编户的合作，而非仅以高额征发维持关中和边郡。

前一七七年前后，文帝弛山泽禁并许可民间铸钱。钱币范铸遗存和《食货志》可以相互提示货币政策的变化，政策适用的范围与实际流通却仍有争议。放宽管制可能减轻中央直接经营的负担，也让诸侯与民间拥有更多财力和铸钱空间；币制质量和地方资源由此成为后来朝廷不得不重新面对的问题。所谓“休养”并非国家从经济中退出，而是它在财政能力有限时改变管制、赋役与信用的组合。

前一六七年废肉刑，改以笞刑等处罚，同样不宜写成司法痛苦已经消失。缇萦故事的传本层次和政策讨论的实际过程需要辨析；法定肉刑被废这一点可靠，而笞刑的死亡风险和地方司法如何执行，不能由诏令本身推出。刑罚的调整说明皇帝可以在法律语言中塑造仁政，也提醒人们，规范与监狱、审讯和地方官实际操作之间始终有距离。

\subsection{边境与封国仍在扩张}

前一六六年匈奴入寇至雍地，长安西北的安全受到直接威胁。入寇兵数不可照录，但和亲并未解除草原机动、边郡驻军与关中补给之间的矛盾。文帝朝必须讨论将领任用和边防，却不能在北方道路、骑兵与粮运条件不足时轻易改写既有关系。帝国能够减租、改刑，并不表示它不必为边境支付代价。

同一时期，贾谊论削藩，文帝又分封诸子为王。前者说明朝廷已意识到王国拥有土地、官属和继承网络所带来的风险；后者则显示刘氏仍要以宗室封国安置皇子并把要地置于本家名义之下。贾谊文字的成篇过程与具体上书年份有争议，诸王各自受封的年月也不完全一致，但这两类行动并不矛盾：中央正试图在宗室依靠和宗室防范之间维持危险的平衡。

\subsection{景帝所继承的不是一张白纸}

前一五七年文帝死，太子启即位，是为景帝。他接手的是减税、法律调整、边防压力和更为密集的同姓王网络。文景时期的政治不能被写成武帝扩张前的一段静止空白；田租、货币、刑罚、边郡和封国已在共同改变帝国的日常运行。削藩最终在前一五四年引发七国之乱，并非某一条建议突然制造的后果，而是这套平衡长期累积的断裂。王国依然存在，中央已不愿再允许它们拥有可以独立调动的军政能力。

\subsection{材料与尺度}

《汉书》〈文帝纪〉〈食货志〉〈刑法志〉、诸侯王表及《史记》相关本纪、列传，构成恢复时期的年序和政策主干；钱范、钱币与边塞材料能补正制度和边防的物质尺度。正史的“文景之治”框架往往以后来的富庶回望早期措施，地方行政簿籍又不足以复原普通家庭的税役经验。本节据此只写诏令、政策争议、边境压力和封国关系怎样互相限制，不以政治赞语代替社会的全貌。
